\documentclass[a4paper]{article}
% Español (México) con XeLaTeX: fontspec para fuentes Unicode, polyglossia
% para la división silábica y los nombres del documento en la variante mexicana.
\usepackage{fontspec}
\usepackage{polyglossia}
\setdefaultlanguage[variant=mexican]{spanish}
% Los nombres chinos van como «pinyin (original)»: xeCJK da a los caracteres
% Han su propia fuente; sin ella XeLaTeX los omite con un aviso.
% FandolSong no cubre algunos caracteres de nombres (U+73FA); WenQuanYi Zen Hei sí.
\usepackage[AutoFallBack=true]{xeCJK}
\setCJKmainfont{FandolSong-Regular.otf}
\setCJKfallbackfamilyfont{\CJKrmdefault}{WenQuanYi Zen Hei}
% polyglossia no traduce los nombres de listings.
\AtBeginDocument{\providecommand{\lstlistingname}{}\renewcommand{\lstlistingname}{Listado}%
  \providecommand{\lstlistlistingname}{}\renewcommand{\lstlistlistingname}{Índice de listados}}
\usepackage{amsmath,amssymb}
\usepackage{graphicx}
\usepackage[margin=2.35cm]{geometry}
\usepackage[most]{tcolorbox}
\usepackage{etoolbox}
\usepackage{listings}
\usepackage{xcolor}
\usepackage{hyperref}
\usepackage{booktabs,longtable,tabularx,array}
\usepackage{subcaption}
\usepackage{float}
\usepackage{enumitem}
\usepackage{microtype}
\usepackage{fancyhdr}
\usepackage{needspace}

\definecolor{kbBack}{HTML}{F2F6FA}
\definecolor{kbFrame}{HTML}{9DB4CC}
\definecolor{kbTitle}{HTML}{52708F}
\definecolor{ibBack}{HTML}{FBF5E7}
\definecolor{ibFrame}{HTML}{C9A961}
\definecolor{ibTitle}{HTML}{7D5F1E}
\definecolor{wbBack}{HTML}{FCF2F0}
\definecolor{wbFrame}{HTML}{D6A096}
\definecolor{wbTitle}{HTML}{9E4F43}
\definecolor{dbBack}{HTML}{F4F7F3}
\definecolor{dbFrame}{HTML}{AEC2AC}
\definecolor{dbTitle}{HTML}{5F7F64}

\tcbset{softbox/.style={
  enhanced,breakable,boxrule=0.8pt,arc=2pt,left=3mm,right=3mm,
  top=2mm,bottom=2mm,coltitle=white,fonttitle=\bfseries,
  toptitle=0.6mm,bottomtitle=0.6mm
}}
\newtcolorbox{knowledgebox}[1]{softbox,colback=kbBack,colframe=kbFrame,colbacktitle=kbTitle,title={#1}}
\newtcolorbox{importantbox}[1]{softbox,colback=ibBack,colframe=ibFrame,colbacktitle=ibTitle,title={#1}}
\newtcolorbox{warningbox}[1]{softbox,colback=wbBack,colframe=wbFrame,colbacktitle=wbTitle,title={#1}}
\newtcolorbox{dialoguebox}[1]{softbox,colback=dbBack,colframe=dbFrame,colbacktitle=dbTitle,title={#1}}

\lstset{
  language=Python,basicstyle=\ttfamily\small,
  keywordstyle=\color{kbTitle}\bfseries,stringstyle=\color{wbTitle},
  commentstyle=\color{dbTitle},breaklines=true,frame=single,numbers=left,
  numberstyle=\tiny\color{gray},captionpos=b,columns=fullflexible,
  keepspaces=true,showstringspaces=false,extendedchars=true
}
\BeforeBeginEnvironment{lstlisting}{\Needspace{12\baselineskip}}

\hypersetup{colorlinks=true,linkcolor=kbTitle,urlcolor=kbTitle,citecolor=wbTitle}
\setlist{itemsep=0.25em,topsep=0.4em}
\setlength{\parindent}{2em}
\setlength{\parskip}{0.25em}
\newcolumntype{L}[1]{>{\raggedright\arraybackslash}p{#1}}

\providecommand{\notetitle}{Notas del curso: [tema del curso]}
\providecommand{\noteauthors}{Elaboradas a partir de los materiales públicos de Stanford CS25}
\providecommand{\notedate}{Versión reescrita en 2026}
\providecommand{\videochannel}{Stanford Online}
\providecommand{\videopublishdate}{2022}
\providecommand{\videoduration}{[duración del video]}
\providecommand{\videourl}{}
\providecommand{\videocoverpath}{cover.jpg}
\providecommand{\coursesite}{https://web.stanford.edu/class/cs25/}
\providecommand{\repourl}{https://github.com/hqhq1025/ai-course-notes}
\providecommand{\skillversion}{wdkns/wdkns-skills@39f1a04}

\newcommand{\figsource}[1]{\par\vspace{-0.3em}{\footnotesize\color{gray}\emph{Fuente: #1}}\par}
\newcommand{\teachervoice}[1]{\begin{knowledgebox}{Nota de clase}#1\end{knowledgebox}}
\newcommand{\term}[1]{\textbf{#1}}

\pagestyle{fancy}
\fancyhf{}
\fancyfoot[C]{\thepage}
\fancyfoot[R]{\footnotesize\color{gray}\href{\repourl}{AI Course Notes}}
\renewcommand{\headrulewidth}{0pt}
\renewcommand{\footrulewidth}{0pt}

\newcommand{\makecscover}{
\begin{titlepage}
\centering
\vspace*{0.7cm}
{\Large Stanford CS25 · Transformers United\par}
\vspace{0.8cm}
{\huge\bfseries \notetitle\par}
\vspace{0.6cm}
{\large \noteauthors\par}
\vspace{0.25cm}
{\large \notedate\par}
\vspace{0.9cm}
\ifdefempty{\videocoverpath}{}{
  \includegraphics[width=0.86\textwidth,height=0.43\textheight,keepaspectratio]{\videocoverpath}\par
}
\vfill
\begin{tcolorbox}[width=0.92\textwidth,colback=black!2!white,colframe=black!45,boxrule=0.7pt,arc=2pt]
\textbf{Autor o canal del video}: \videochannel\par
\textbf{Fecha de publicación}: \videopublishdate\par
\textbf{Duración del video}: \videoduration\par
\textbf{Sitio del curso}: \href{\coursesite}{\nolinkurl{\coursesite}}\par
\ifdefempty{\videourl}{}{\textbf{Enlace del video}: \href{\videourl}{\nolinkurl{\videourl}}\par}
\textbf{Normas de elaboración}: \skillversion\ + las normas source-first / teacher-voice / visual-QA de este repositorio
\end{tcolorbox}
\end{titlepage}
\tableofcontents
\newpage
}

\usepackage{multicol}

\renewcommand{\notetitle}{Lecture 31: De los grandes modelos de lenguaje a los grandes modelos multimodales}
\renewcommand{\noteauthors}{Elaborado a partir de la clase de Ming Ding, las slides recuperadas del video oficial y los subtítulos manuales}
\renewcommand{\notedate}{CS25 V4 · Versión reescrita source-first de 2026}
\renewcommand{\videochannel}{Stanford Online}
\renewcommand{\videopublishdate}{Clase: 2024-05-09; publicación: 2024-05-30}
\renewcommand{\videoduration}{1:20:03}
\renewcommand{\videourl}{https://www.youtube.com/watch?v=cYfKQ6YG9Qo}
\renewcommand{\videocoverpath}{cover.jpg}

\newcommand{\lecturefigure}[3]{%
\begin{figure}[H]
  \centering
  \includegraphics[width=0.97\textwidth,height=0.49\textheight,keepaspectratio]{slides-images/#1}
  \caption{#2}
  \figsource{#3}
\end{figure}
}

\let\standardtableofcontents\tableofcontents
\renewcommand{\tableofcontents}{%
  \begingroup
  \small
  \setlength{\parskip}{0pt}
  \standardtableofcontents
  \endgroup
}

\begin{document}
\makecscover

\section{Ubicación de la clase: la multimodalidad no consiste en añadir un módulo de visión por fuera del LLM}

Esta clase abarca mucho terreno. La primera parte repasa tres «momentos» de los modelos de lenguaje: BERT, GPT-3 y ChatGPT. La parte intermedia trata la arquitectura Transformer, los sistemas de entrenamiento, la alignment y la ingeniería de datos. Solo la última parte llega a BLIP-2, LLaVA, CogVLM, CogAgent, la difusión de imágenes y la generación de video. Este orden no es una digresión, sino la respuesta a una pregunta que recorre toda la clase: \textbf{cuando las capacidades del modelo se extienden de text a image, GUI y video, ¿qué regularidades provenientes de los LLM siguen siendo válidas y cuáles exigen cambiar de representación, de función objetivo o de ruta de sistema?}

\subsection{Corrección de fuentes y límites de la clase}

El borrador anterior fechaba la clase el 4 de abril de 2026 y añadía mucho contenido que no existe en el curso, como monitoring dashboard, drift detection, automatic rollback y fairness governance. Ahora esta clase toma como eje la página oficial del curso CS25 V4, la sesión del 9 de mayo de 2024, la grabación oficial de Stanford Online \texttt{cYfKQ6YG9Qo}, los subtítulos humanos oficiales en \texttt{en-US} y 31 teaching states recuperados de la grabación. Ni la página oficial ni la descripción del video ofrecen un slide deck independiente que pueda verificarse; por eso, la pantalla compartida en la grabación es la source of truth visual.

\lecturefigure{slide-01-title.jpg}{Título oficial de la clase: From Large Language Models to Large Multi-modality Models}{Grabación oficial de Stanford Online 00:02:31}

El término ``Large Multi-modality Models'' del título es el que se usaba en la clase de 2024. No significa que todas las modality ya se hayan resuelto de manera unificada; subraya que la investigación sobre modelos ya pasó de una sola interfaz de lenguaje a la comprensión visual, la generación de imágenes, las GUI action y el video. CogVLM, CogAgent, CogView y CogVideo, del equipo del ponente, son casos de estudio y no la única conclusión del curso.

\begin{warningbox}{Hay que conservar el límite temporal}
Los benchmark, las capacidades de producto, las cifras de descargas y las predicciones «a uno o dos años» que se muestran en esta clase son una instantánea de la sesión del 2024-05-09. Las versiones posteriores de los modelos, los cambios en las tablas de clasificación o los resultados de la industria no deben proyectarse hacia atrás como hechos de la clase; si se usan como ampliación, deben marcarse explícitamente como post-lecture evidence.
\end{warningbox}

\subsection{Hoja de ruta en cuatro partes y preguntas de lectura}

El curso primero pregunta why, luego how, después what y por último where. El lector no debe ver las cuatro partes como una revisión de temas paralelos. El objective y el scaling de los LLM explican por qué la ingeniería de sistemas cobró importancia. Los sistemas y los datos explican por qué los VLM pueden reutilizar con rapidez la base de lenguaje. A su vez, las carencias de la representación visual explican por qué la generación pasó de la autoregression a la diffusion.

\lecturefigure{slide-02-outline.jpg}{Hoja de ruta del curso: historia, técnicas de entrenamiento, modelos de visión y lenguaje y líneas de investigación}{Grabación oficial de Stanford Online 00:03:29}

\begin{importantbox}{Cuatro preguntas para leer toda la clase}
\begin{enumerate}
  \item \textbf{Objective:} ¿las distintas modality deben compartir un autoregressive objective o usar un objective especializado de diffusion / contrastive?
  \item \textbf{Architecture:} ¿qué problema de interfaz resuelven, cada uno, la capa de proyección, el Q-Former, los vision experts y la cross attention?
  \item \textbf{Systems:} ¿cómo cambian el límite de los problemas entrenables ZeRO (Zero Redundancy Optimizer, que reparte los optimizer/gradient/parameter states entre los data-parallel ranks), el parallelism, el long context y los token de alta resolución?
  \item \textbf{Data:} cuando los web data generales se acercan a la saturación, ¿cómo se convierte el compute en datos de mayor calidad y verificables?
\end{enumerate}
\end{importantbox}

\teachervoice{00:02:31--00:04:35. Ming Ding advierte al público que la clase recorrerá varios subcampos que quizá no conozca bien. Su objetivo no es hacer un survey completo de cada modelo, sino explicar, siguiendo la cronología de la comunidad Transformer, por qué se desplazó el foco de la investigación.}

\subsection{Registro de evidencia: mecanismos, resultados, juicios y predicciones}

Una misma slide puede contener a la vez un diagrama de arquitectura de un artículo, un juicio expresado en clase y un resultado de producto. Para evitar que «lo dice la figura» se lea como un hecho general, esta clase usa un registro de evidencia en cuatro niveles.

\begin{knowledgebox}{Cómo leer la evidencia en cuatro niveles}
\begin{tabularx}{\textwidth}{L{0.18\textwidth} L{0.31\textwidth} X}
\toprule
Nivel & Ejemplos & Lectura correcta \\
\midrule
Mecanismo estándar & cross entropy, attention, DDPM, DPO & Dar fórmulas, variables y supuestos \\
Sistema concreto & ZeRO, Megatron, CogVLM, MM-DiT & Explicar interfaces, rutas de recursos y versiones \\
Resultado de clase & benchmark table, web-agent trace, cifras de velocidad & Solo respaldan la configuración y el protocol mostrados \\
Juicio/predicción del ponente & la loss determina las capacidades; el video será un tema central & Conservar la fecha, el tono y el margen para contraejemplos \\
\bottomrule
\end{tabularx}
\end{knowledgebox}

\subsection{Resumen de la sección}

Esta es una clase sobre «trasladar las regularidades de los LLM a la multimodality», no un manual de producto de la serie Cog. El lector debe seguir a la vez el objective, la architecture, los systems y los data, y distinguir siempre entre mecanismos estándar, experimentos de clase, juicios del ponente y predicciones a futuro.
\section{Tres LLM moments: función objetivo, escala y alineación}

El ponente condensa cinco años de historia de los modelos de lenguaje en tres momentos: BERT, GPT-3 y ChatGPT. Cada uno cambió la forma en que los investigadores entendían la self-supervision, la scale y la task adaptation. Lo importante no es memorizar los años, sino observar cómo el cuello de botella de la investigación pasó de «diseñar mejores tareas» a «asignar el cómputo», y después a «cambiar el comportamiento con pocos datos de alta calidad».

\subsection{BERT moment: se creía que la comprensión y la generación requerían objetivos distintos}

Entre 2018 y 2021, la división de trabajo habitual era esta: el masked language modeling (MLM) se encargaba de la comprensión, el autoregressive modeling de la generación, y T5 usaba encoder--decoder para cubrir ambas. El blank infilling de GLM, en cambio, convierte el interior de cada span enmascarado en generación autorregresiva, con la intención de cubrir tareas distintas con un solo decoder-style objective.

\lecturefigure{slide-03-language-modeling.jpg}{El momento BERT: MLM, autoregression, T5 y GLM desde la perspectiva de la tarea}{Grabación oficial de Stanford Online 00:07:00}

Sea la secuencia original $x=(x_1,\ldots,x_n)$ y sea $M$ el conjunto de spans enmascarados. El núcleo de GLM puede escribirse como una autorregresión condicional sobre los token enmascarados:
\[
\mathcal{L}_{\text{GLM}}
=-\sum_{m\in M}\sum_{j=1}^{|m|}
\log p_\theta\!\left(x_{m,j}\mid x_{\setminus M},x_{m,<j}\right).
\]
Aquí, $x_{\setminus M}$ es el contexto visible y $x_{m,<j}$ son los token ya generados dentro del span actual. Si el enmascaramiento cubre casi toda la secuencia, el objetivo se aproxima al de GPT; si se enmascaran spans pequeños, se conserva el contexto visible bidireccional.

\begin{knowledgebox}{La verdadera función del objective}
El objective de preentrenamiento determina qué condiciones ve el modelo durante el entrenamiento, qué variables predice y cómo se propaga el error. No equivale directamente a la capacidad de «comprensión» o de «generación»: el fine-tuning posterior, el prompt format, la data mixture y el inference procedure también modifican el comportamiento final.
\end{knowledgebox}

\teachervoice{00:05:00--00:08:55: el ponente sitúa GLM en una historia en la que «las creencias de la comunidad cambian con rapidez». En su momento, todos se esforzaban por inventar self-supervised tasks; después se descubrió que la task adaptation podía ser muy barata, y la valoración de qué investigación aporta cambia según la escala y las condiciones de los datos.}

\subsection{GPT-3 moment: la scaling law se convierte en una herramienta para asignar el presupuesto}

El segundo momento no es el lema de que «más parámetros es mejor», sino el uso de curvas ajustables que relacionan la loss con los parámetros, los datos y el cómputo. Una forma empírica habitual es
\[
L(N,D)\approx L_\infty+aN^{-\alpha}+bD^{-\beta},
\]
donde $N$ es el número de parámetros, $D$ es el número de token de entrenamiento, $L_\infty$ es la loss irreducible y $a,b,\alpha,\beta$ se ajustan experimentalmente. Esto permite a un equipo comparar, con un compute $C$ fijo, el rendimiento marginal de «agregar parámetros» frente al de «agregar token».

\lecturefigure{slide-04-scaling-law.jpg}{El momento GPT-3: usar scaling curves para asignar parámetros, datos y cómputo}{Grabación oficial de Stanford Online 00:09:19}

Si los FLOPs principales de un entrenamiento de un dense Transformer se aproximan como
\[
C\approx 6ND,
\]
entonces, con un presupuesto dado, $N$ y $D$ no pueden crecer de manera ilimitada al mismo tiempo. El $6$ es una constante aproximada de uso común; el costo real también depende de la sequence length, la attention, la activation recomputation, la parallel efficiency y los kernels.

\begin{warningbox}{La scaling law no es una garantía de rendimiento}
La curva ajustada solo es válida dentro de la familia de modelos, la distribución de datos, el tokenizer, el optimizer y el intervalo de entrenamiento observados. Si baja la calidad de los datos, aumenta la tasa de repetición, se trata de tareas fuera de distribución o empeora la eficiencia del sistema, aumentar el compute nominal no garantiza una mejora con el mismo exponente.
\end{warningbox}

\teachervoice{00:09:19--00:10:00: el ponente describe la scaling law como una herramienta de presupuesto de ingeniería: cuando el jefe asigna cuatro veces más recursos, hay que decidir si se destinan a parámetros o a token, en lugar de limitarse a anunciar que «el modelo es más grande».}

\subsection{ChatGPT moment: la task adaptation es barata, pero ¿de dónde viene el conocimiento?}

La tercera slide reúne la mejora en preferencia humana de InstructGPT y la gráfica empírica según la cual «el rendimiento posterior lo determina la pretraining loss». Lo primero muestra que el alignment puede cambiar de forma significativa la evaluación humana; lo segundo intenta mostrar que, con la misma loss de preentrenamiento, un modelo más pequeño pero entrenado de manera más completa puede alcanzar a uno más grande pero under-trained.

\lecturefigure{slide-05-alignment-pretraining-loss.jpg}{El momento ChatGPT: el alignment cambia el comportamiento; la pretraining loss sostiene las capacidades básicas}{Grabación oficial de Stanford Online 00:13:05}

La relación entre la entropía cruzada y la perplexity es
\[
H=-\frac{1}{T}\sum_{t=1}^{T}\log p_\theta(x_t\mid x_{<t}),
\qquad
\operatorname{PPL}=e^H.
\]
La perplexity puede entenderse de forma intuitiva como «el número efectivo de opciones que tiene el modelo ante cada token», pero la PPL no puede compararse directamente entre tokenizers, corpus y domains distintos.

\begin{warningbox}{«Misma loss, misma capacidad» no es un teorema general}
Lo que se plantea en clase es un juicio empírico fuerte. Aunque la pretraining loss promedio sea similar, la architecture, el tokenizer, la data mixture, la optimization trajectory, la context length y el evaluation format pueden seguir produciendo diferencias de capacidad. Los trabajos posteriores a septiembre de 2024 también presentan contraejemplos de same-loss, different-downstream; por eso, esta clase solo conserva la intuición sobre el presupuesto y el over-training.
\end{warningbox}

\teachervoice{00:10:00--00:13:18: el ponente reinterpreta explícitamente la llamada emergent ability como un fenómeno de la curva de loss. Esta es su postura de investigación y no debe reescribirse como una conclusión causal que la comunidad ya haya demostrado.}

\subsection{Resumen de la sección}

El BERT moment se centra en el objetivo de entrenamiento; el GPT-3 moment, en la escala y el presupuesto; el ChatGPT moment, en la adaptación del comportamiento a bajo costo. Los tres, en conjunto, transformaron la pregunta de investigación: pasó de «qué tarea es la más inteligente» a «con qué datos y qué cómputo se obtiene el conocimiento, y cómo se moldea el comportamiento mediante retroalimentación de alta calidad».
\section{Arquitectura Transformer y sistemas de entrenamiento: los recursos marcan el límite de las capacidades}

Antes de que aparecieran los modelos multimodales, la base de lenguaje tenía que entrenarse de forma estable. La tesis central del ponente es que muchas de las mejoras clave de los LLM modernos no son redes completamente nuevas, sino ajustes continuos al decoder path, la normalization, la KV memory, el parameter routing y la distributed execution. Entender estos términos es requisito para entender después los VLM de high-resolution y el video scaling.

\subsection{Términos clave de la recipe decoder-only moderna}

La sección anterior explicó cómo el objective, la scale y el alignment desplazan el cuello de botella de la investigación; esta sección va un paso más allá y pregunta: una vez que el LLM decoder-only se convierte en la base común, ¿qué partes del grafo de cómputo modificaron realmente los equipos de ingeniería, y qué ganó cada cambio en estabilidad, capacidad o inference memory? La siguiente slide reúne en una sola página las architecture updates más comunes. Al leer la figura no basta con contar nombres: hay que distinguir si cada una actúa sobre el residual path, la position, el KV sharing, el MLP o el routing.

\lecturefigure{slide-06-transformer-architecture.jpg}{Actualizaciones comunes del Transformer moderno: decoder-only, pre-norm, RoPE, GQA, GLU y MoE}{Grabación oficial de Stanford Online 00:16:10}

La causal self-attention estándar es
\[
Q=XW_Q,\quad K=XW_K,\quad V=XW_V,\qquad
\operatorname{Attn}(X)=\operatorname{softmax}\!\left(\frac{QK^\top}{\sqrt{d_h}}+M\right)V,
\]
donde $M$ es la causal mask. GQA (Grouped-Query Attention) hace que varios grupos de query heads compartan menos key/value heads; su efecto principal es reducir la KV cache, no reducir todos los FLOPs de attention.

\begin{knowledgebox}{Glosario denso: cada término resuelve un cuello de botella distinto}
\begin{tabularx}{\textwidth}{L{0.17\textwidth} L{0.30\textwidth} X}
\toprule
Término & Mecanismo central & Relación con el curso \\
\midrule
Decoder-only & Un solo causal generation path & Facilita escalar el preentrenamiento y unificar la interfaz \\
Pre-norm & LayerNorm antes de cada subcapa & Mejora la estabilidad del entrenamiento con residuales profundos \\
RoPE & Codifica la posición relativa rotando $Q,K$ & Da estructura posicional para contexto largo y extrapolación \\
GQA & Varias query heads comparten KV heads & Reduce la KV cache en inference \\
SwiGLU/GLU & Sustituye la activación del MLP común por una multiplicación con compuerta & Cambia la capacidad y la optimización del feed-forward \\
MoE & Cada token activa solo una parte de los experts & Aumenta los parámetros totales sin descontrolar el active compute \\
\bottomrule
\end{tabularx}
\end{knowledgebox}

\teachervoice{00:13:18--00:17:45, el ponente enfatiza que la mayoría de estas mejoras se acumularon durante mucho tiempo y no surgieron de golpe como un «nuevo Transformer». Bromea con que hay muy poca architecture innovation, pero eso no significa que cada cambio sea irrelevante.}

\subsection{DeepSpeed y ZeRO: primero calcular cuánta memoria de GPU ocupa cada parámetro}

La primera lección de los sistemas de entrenamiento no es memorizar ZeRO-1/2/3, sino hacer resource accounting. Con AdamW en precisión mixta, cada parámetro suele requerir el parameter y el gradient en FP16/BF16, además del master weight, el first moment $m$ y el second moment $v$ en FP32. El número aproximado de bytes es
\[
B_{\text{per-param}}\approx 2_{\theta}+2_g+4_{\theta_{32}}+4_m+4_v=16\ \text{bytes}.
\]
Con mil millones de parámetros, solo estos states ocupan unos 16 GB, sin contar activation, temporary buffers, fragmentation ni el communication workspace.

\lecturefigure{slide-07-deepspeed.jpg}{DeepSpeed: optimizer states, activation checkpointing y fragmentación con ZeRO}{Grabación oficial de Stanford Online 00:19:30}

\begin{importantbox}{Primera definición: sharding, optimizer state y activation checkpointing}
\textbf{Sharding (fragmentación)} consiste en repartir entre distintos data-parallel ranks un tipo de estado que antes se replicaba en cada GPU; el \textbf{optimizer state} son los $m,v$ de Adam/AdamW y los habituales master weights en FP32; el \textbf{activation checkpointing} guarda solo una parte de las forward activations y recalcula el resto durante el backward, lo cual no es lo mismo que el model checkpointing, que guarda el archivo del modelo.
\end{importantbox}

ZeRO-1 fragmenta el optimizer state, ZeRO-2 fragmenta además los gradients y ZeRO-3 fragmenta también los parameters. Si el data-parallel world size es $p$, en el caso ideal la memoria ocupada por los estados se reduce aproximadamente en $1/p$, pero a cambio aparecen collectives (primitivas de comunicación entre varias GPU) como all-gather y reduce-scatter.

\begin{lstlisting}[caption={Lista de cálculo de memoria de GPU para entrenamiento en precisión mixta}]
bytes = parameter_low_precision
bytes += gradient_low_precision
bytes += master_parameter_fp32
bytes += adam_first_moment_fp32
bytes += adam_second_moment_fp32
bytes += saved_activations_or_recompute_budget
bytes += communication_and_kernel_workspace
report(bytes_per_rank, sharding_stage, world_size)
\end{lstlisting}

\teachervoice{00:17:45--00:21:13, el ponente señala en particular que la mayor parte de la memoria de GPU suele provenir de los Adam states, y que el costo del activation checkpointing es el recálculo en el backward. Presentar el «ahorro de memoria» como una optimización gratuita omite un hecho del sistema.}

\subsection{Megatron: qué divide la tensor parallelism y qué divide la pipeline parallelism}

ZeRO fragmenta los estados a lo largo de los data-parallel ranks; la tensor parallelism (TP) de Megatron divide las hidden dimensions/heads dentro de una capa, y la pipeline parallelism (PP) divide en stages a lo largo de la profundidad de las capas. Ambas actúan sobre dimensiones distintas y también tienen patrones de comunicación distintos.

\lecturefigure{slide-08-megatron.jpg}{Megatron: cómo dividen el modelo la tensor parallel y la pipeline parallel}{Grabación oficial de Stanford Online 00:21:13}

Si la capa lineal $Y=XW$ se divide por columnas, $W=[W_1,\ldots,W_p]$, cada rank calcula $Y_i=XW_i$ y después puede requerirse un all-gather; si se divide por filas, los resultados locales necesitan un all-reduce para sumarse. La utilización de la PP depende del número de micro-batches $m$ y del número de stages $p$; la bubble fraction ingenua puede escribirse aproximadamente como
\[
\rho_{\text{bubble}}\approx \frac{p-1}{m+p-1}.
\]
Las planificaciones de tipo interleaving, 1F1B y ZeroBubble reducen el tiempo ocioso, pero aumentan la complejidad de la programación y las restricciones de comunicación.

\begin{knowledgebox}{Las collectives no son una «sincronización» difusa}
\textbf{All-reduce} agrega los datos y devuelve el resultado a cada rank; \textbf{reduce-scatter} agrega y cada rank conserva solo su fragmento; \textbf{all-gather} reúne todos los fragmentos; \textbf{broadcast} envía desde un rank a todos los ranks. Las diferencias de rendimiento entre distintos esquemas de sharding/parallelism suelen deberse a la frecuencia de estas operaciones, al tamaño de los mensajes y a su correspondencia con la topología.
\end{knowledgebox}

\teachervoice{00:20:00--00:22:50, el ponente resume el entrenamiento de modelos muy grandes como engineering work y, con la broma de que «el MLP no importa, lo que importa es ML6», enfatiza que sin systems knowledge los FLOPs teóricos de la arquitectura no se convierten en un entrenamiento efectivo.}

\subsection{Long context: la full attention sigue pagando una complejidad cuadrática}

La slide de contexto largo compara la estructura explícita de retrieval/rehearsal del antiguo CogLTX con la full-attention context parallelism de 2024. Lossless significa que no se modifica el attention pattern mediante sparse attention, no que el cálculo no tenga costo.

\lecturefigure{slide-09-long-context.jpg}{Lossless 128K: context parallel, Ring Attention y Ulysses}{Grabación oficial de Stanford Online 00:24:27}

Para una sequence length $n$ y una head dimension $d_h$, el cálculo principal de la score matrix completa es
\[
\mathcal{O}(n^2d_h),
\]
y su almacenamiento explícito en una sola GPU también puede llegar a $\mathcal{O}(n^2)$. La context parallelism reparte la sequence entre $p$ ranks y reduce la activation local a cerca de $\mathcal{O}(nd/p)$, pero obliga a rotar o intercambiar bloques de $K,V$ y a manejar el causal load imbalance.

\begin{warningbox}{El costo oculto del contexto largo es el prefill}
En la sesión de preguntas y respuestas, el ponente divide la inference en prefill y decode. Con documentos largos a menudo solo se genera una respuesta corta, pero el prefill tiene que leer y procesar todo el context, lo que puede tardar desde unos segundos hasta un minuto. Una ventana más larga no es «memory gratuita»: también aumenta la latency, la KV cache, la communication y la interferencia de información irrelevante.
\end{warningbox}

\teachervoice{01:08:10--01:10:30, el ponente acepta que en ciertos escenarios de comprensión de documentos se pague un prefill más largo a cambio de una respuesta corta, pero no afirma que todas las aplicaciones interactivas puedan tolerar esa latency.}

\subsection{Resumen de la sección}

La recipe moderna de los LLM incluye a la vez architecture updates y distributed systems. ZeRO divide los estados, la TP divide los tensores dentro de la capa, la PP divide las capas y la context parallelism divide la secuencia; cada reducción de memoria por GPU se paga con collectives, recálculo, bubble o latency. Estas rutas de sistema volverán a aparecer en las high-resolution images y en el video.
\section{Alignment e ingeniería de datos: adaptar la conducta no es generar conocimiento de la nada}

Una vez entrenada la base del lenguaje, SFT y preference optimization determinan cómo responde el modelo. El ponente agrupa esta parte con data engineering porque la variable decisiva de alignment no es el nombre del algoritmo, sino la calidad, la procedencia y la cobertura de las instruction, las answer y los preference pair.

\subsection{SFT: las respuestas de alta calidad requieren domain experts}

SFT (Supervised Fine-Tuning) sigue siendo la cross entropy teacher-forced habitual. Dado un par instruction--response $(x,y)$, el objetivo es
\[
\mathcal{L}_{\text{SFT}}=-\sum_{t=1}^{|y|}\log \pi_\theta(y_t\mid x,y_{<t}).
\]
La dificultad está en los datos: si se quiere que el modelo escriba explicaciones de código confiables, quien anota debe ser capaz de escribir explicaciones de código confiables, y no limitarse a resolver preguntas de opción múltiple de bajo costo.

\lecturefigure{slide-10-sft.jpg}{SFT: expert annotation, distillation y el problema weak-to-strong}{Grabación oficial de Stanford Online 00:27:01}

\begin{knowledgebox}{Los dos límites de los teacher-generated data}
Un modelo más fuerte puede generar instruction data con rapidez, pero el estudiante hereda los sesgos, el estilo y los puntos ciegos del teacher; si el objetivo es competir de forma independiente, depender de un teacher cerrado trae además problemas de licencias, reproducibilidad y data provenance. Weak-to-Strong estudia si, bajo supervisión débil, se puede superar al supervisor; no equivale a «copiar las teacher outputs basta para superar necesariamente al teacher».
\end{knowledgebox}

\teachervoice{00:25:00--00:28:24, el ponente insiste en que SFT no es crowdsourcing ordinario. Las respuestas de alta calidad requieren personas que de verdad conozcan el domain, y esto importa más que reducir la cantidad de datos a una cifra.}

\subsection{RLHF y DPO: el algoritmo se simplifica, pero el supuesto de preferencias no desaparece}

SFT enseña al modelo a imitar la respuesta objetivo, pero no compara de forma explícita dos salidas que «parecen respuestas» por igual. Preference optimization responde entonces la siguiente pregunta: cuando chosen y rejected solo difieren en utilidad, estilo o seguridad, cómo convertir un juicio por pares en una policy update. La diferencia entre RLHF basado en PPO y DPO está en la ruta de optimización, no en si hace falta definir las preferencias.

RLHF al estilo PPO suele entrenar primero un reward model $r_\phi(x,y)$ y después optimizar la policy con KL regularization:
\[
\max_\theta\ \mathbb{E}_{y\sim\pi_\theta(\cdot\mid x)}[r_\phi(x,y)]
-\beta D_{\mathrm{KL}}\!\left(\pi_\theta\,\|\,\pi_{\text{ref}}\right).
\]
DPO optimiza directamente con el pair chosen/rejected $(y^+,y^-)$
\[
\mathcal{L}_{\text{DPO}}
=-\log\sigma\!\left(\beta\left[
\log\frac{\pi_\theta(y^+\mid x)}{\pi_{\text{ref}}(y^+\mid x)}
-\log\frac{\pi_\theta(y^-\mid x)}{\pi_{\text{ref}}(y^-\mid x)}
\right]\right).
\]

\lecturefigure{slide-11-rlhf.jpg}{RLHF y DPO: reward model, la dificultad de PPO y el pairwise objective}{Grabación oficial de Stanford Online 00:28:24}

\begin{warningbox}{DPO no elimina el reward modeling}
DPO no entrena de forma explícita un reward model independiente, pero los preference pairs siguen definiendo de manera implícita «qué es mejor». Que el pair sea on-policy o no, que los annotator sean consistentes, que los prompt cubran el uso real y que chosen/rejected difieran solo en estilo: todo ello modifica la reward geometry que se aprende.
\end{warningbox}

\teachervoice{00:28:24--00:30:00, el ponente considera que PPO puede ser muy potente cuando el reward model es lo bastante bueno, pero su implementación es difícil; presenta DPO como una alternativa más fácil de usar y, al mismo tiempo, advierte de palabra sobre la pair quality y el carácter on-policy.}

\subsection{Data, algorithm y architecture pueden codificar un bias similar}

Esta slide contiene la metodología de investigación más importante de toda la clase. El multi-hop question answering de CogQA llegó a usar una graph neural network, y otra línea usó MCTS; cuando un GPT de long-context y chain-of-thought pueden colocar de una sola vez en el context todos los documents relevantes, la misma tarea puede reformularse como un problema a nivel de data/context.

\lecturefigure{slide-12-data-algorithm-architecture.jpg}{Intercambiabilidad entre datos, algoritmo y arquitectura: de CogQA al razonamiento con long-context}{Grabación oficial de Stanford Online 00:34:16}

\begin{importantbox}{Tres niveles para implementar el mismo sesgo inductivo}
\begin{enumerate}
  \item \textbf{Architecture:} incorporar graph structure, memory o routing en la estructura de la red.
  \item \textbf{Algorithm:} organizar el cómputo en inference mediante search, MCTS, retrieval o un verifier.
  \item \textbf{Data/context:} colocar la evidencia intermedia, las rutas fallidas o los documents completos en los ejemplos de entrenamiento y en el contexto, para que un modelo general aprenda la procedure.
\end{enumerate}
Cuanto más cerca del nivel de los datos, más general suele ser la solución, pero exige un modelo más fuerte, un context más largo y ejemplos de mayor calidad.
\end{importantbox}

\begin{lstlisting}[caption={Marco de decisión para elegir entre data, algorithm o architecture}]
if behavior_is_narrow_and_examples_are_available:
    encode_with_data_and_evaluation()
elif inference_needs_verifiable_search:
    add_algorithmic_search_or_tools()
elif bottleneck_is_general_and_repeats_across_tasks:
    consider_architecture_change()
always_measure(compute, latency, coverage, failure_modes)
\end{lstlisting}

\teachervoice{00:30:00--00:34:20, el ponente defiende el data work: cleaning, filtering y synthesizing no son «trabajo de bajo nivel», sino la ingeniería de investigación más central de las empresas de modelos actuales.}

\teachervoice{01:10:30--01:12:30, en la sesión de Q\&A añadió un matiz: muchas conductas concretas pueden resolverse con data, pero si alguien encuentra una architecture update que mejore de forma general la capacidad de ajuste, sigue teniendo mucho valor.}

\subsection{Resumen de la sección}

SFT, RLHF y DPO no pueden discutirse al margen de la calidad de los datos. En términos más generales, un mismo inductive bias puede incorporarse en architecture, algorithm o data/context; el nivel más sencillo no es necesariamente el más barato, porque puede exigir más compute, un context más largo y una evaluation más confiable.
\section{Interfaces entre visión y lenguaje: de Q-Former a la proyección lineal}

Al conectar image con un LLM, la primera pregunta no es «¿el modelo entendió?», sino cómo entran las características continuas del encoder visual en el token space del language model. BLIP-2 y LLaVA ofrecen dos formas típicas de hacer ese puente: uno usa queries aprendibles y un Q-Former para comprimir y alinear; el otro proyecta linealmente y de forma directa los token visuales.

\subsection{BLIP-2: consultar el espacio visual congelado con un Q-Former}

BLIP-2 coloca un Q-Former entre un frozen image encoder y un frozen LLM. Sean las características visuales $V\in\mathbb{R}^{m\times d_v}$ y las queries aprendibles $Q_0\in\mathbb{R}^{q\times d}$; la cross attention es
\[
Q'=\operatorname{softmax}\!\left(\frac{(Q_0W_Q)(VW_K)^\top}{\sqrt d}\right)VW_V.
\]
Cuando $q\ll m$, el Q-Former funciona a la vez como cuello de botella de información y como alineación entre espacios; después envía al LLM un número reducido de query outputs.

\lecturefigure{slide-13-blip2.jpg}{BLIP-2: el Q-Former sirve de puente entre un frozen image encoder y un frozen LLM}{Grabación oficial de Stanford Online 00:36:18}

\begin{knowledgebox}{El Q-Former resuelve dos problemas distintos}
El primero es la \textbf{compression}: extraer un número fijo de queries a partir de una gran cantidad de image patches. El segundo es la \textbf{alignment}: llevar el feature space de tipo CLIP al embedding space que el LLM puede consumir. No se trata de que el LLM vuelva a aprender por sí mismo todo el backbone visual.
\end{knowledgebox}

\teachervoice{00:35:00--00:36:55, el ponente describe BLIP-2 como pionero en conectar de forma eficaz CLIP con un LLM; el énfasis está en que «entre dos espacios congelados hace falta un puente entrenable».}

\subsection{LLaVA: por qué una projection simple se volvió un baseline sólido}

LLaVA usa un vision encoder $f_v$ y una projection $W_p$:
\[
H_v=f_v(I),\qquad Z_v=H_vW_p,
\]
y luego concatena $Z_v$ con los text embeddings como entrada del LLM. Reduce los módulos adicionales y facilita el instruction tuning end-to-end, por lo que pronto se convirtió en un baseline habitual de los VLM.

\lecturefigure{slide-14-llava.jpg}{LLaVA: alinear las características visuales con el embedding de lenguaje mediante una projection}{Grabación oficial de Stanford Online 00:36:55}

\begin{warningbox}{Que la linear projection sea simple no significa que no pierda información}
La projection solo garantiza que las dimensiones coincidan. La resolución del encoder visual, el patch size, el pretraining objective, la projector capacity y los instruction data determinan qué detalles llegan al LLM; el texto pequeño, el OCR, el grounding y los GUI icons suelen ser los primeros en revelar el cuello de botella.
\end{warningbox}

\subsection{Taxonomy de interfaces: compresión, alineación, conservación y alta resolución}

Antes de pasar a la serie Cog, todos los VLM bridge pueden compararse con cuatro preguntas: ¿cuántos token visuales se conservan? ¿Comparten parámetros la visión y el lenguaje? ¿Se daña la capacidad de lenguaje original? ¿Dónde se paga el cost de la entrada de alta resolución?

\begin{knowledgebox}{Tabla comparativa de VLM bridge}
\begin{tabularx}{\textwidth}{L{0.16\textwidth} L{0.22\textwidth} L{0.23\textwidth} X}
\toprule
Método & Interfaz visual & Ruta de parámetros & Compromiso principal \\
\midrule
BLIP-2 & Q-Former queries & bridge entrenable, backbones congelables & compresión fuerte, pero el query bottleneck es evidente \\
LLaVA & linear/MLP projector & los token visuales entran en la LLM path & línea base simple y sólida; costo alto en alta resolución \\
CogVLM & vision experts & los image/text token pasan en parte por parámetros distintos & intenta conservar el comportamiento de lenguaje; más parámetros \\
CogAgent & high-res cross attention & desvío ligero de alta resolución & orientado a GUI/OCR; sistema más complejo \\
Vary & múltiples vocabularios visuales/feature ensemble & fusión de características de varios encoder & OCR más fuerte; aumentan los costos de interfaz y entrenamiento \\
GLM-4V slide & stride convolution & mixed training tras la compresión & estructura simple; los resultados dependen de los datos y del protocol \\
\bottomrule
\end{tabularx}
\end{knowledgebox}

\subsection{Resumen de la sección}

En esencia, la interfaz multimodal es una cuestión de information bottleneck y parameter sharing. El Q-Former consulta y comprime de forma explícita; LLaVA proyecta directamente. Los modelos posteriores se diversifican en torno a «conservar la capacidad de lenguaje» y «procesar los detalles de alta resolución».
\section{CogVLM, CogAgent y GLM-4V: separar las rutas de parámetros o comprimir la entrada}

Esta sección analiza en conjunto los modelos de comprensión visual que presentó el equipo del ponente. Lo importante no es comparar nombres de modelos, sino observar tres cuellos de botella: si la image alignment sobrescribe la capacidad lingüística, cómo controla una high-resolution screenshot el token cost y si los benchmarks de OCR/grounding pueden representar a un GUI agent real.

\subsection{CogVLM: los vision experts evitan sobrescribir el language path}

La motivación de CogVLM es la siguiente: si todos los image-text alignment gradients actualizan los LLM parameters originales, el comportamiento lingüístico básico puede quedar sobrescrito. El modelo agrega experts independientes de attention/MLP para los tokens visuales y conserva la ruta original de los text tokens. Puede abstraerse como
\[
h'_{i}=h_i+\begin{cases}
F_{\text{vision}}(h_i), & i\in\mathcal{I}_{\text{image}},\\
F_{\text{text}}(h_i), & i\in\mathcal{I}_{\text{text}}.
\end{cases}
\]

\lecturefigure{slide-15-cogvlm-architecture.jpg}{CogVLM: los vision experts separan las rutas de parámetros visuales y lingüísticos}{Grabación oficial de Stanford Online 00:38:44}

\begin{importantbox}{Cómo debe verificarse que se «conserva el comportamiento lingüístico»}
Conservar el text expert en la estructura solo aporta una hipótesis sobre el mecanismo. Una verificación real exige comparar, antes y después del multimodal tuning, la language-only perplexity, el instruction following, el reasoning, la safety y las domain tasks; no se puede afirmar que no hay olvido solo porque los parámetros no se comparten por completo.
\end{importantbox}

\teachervoice{00:37:25--00:39:27, el ponente formula explícitamente el objetivo de diseño de CogVLM como keep all language behavior. Las notas conservan ese objetivo, pero no lo elevan a una conclusión que no ha pasado por pruebas completas de regresión lingüística.}

\subsection{Página de resultados de CogVLM: ante el gráfico de radar, primero el protocol y después el área}

La slide de resultados resume image captioning, grounding, VQA y VLM benchmarks generales, y ofrece los enlaces de código abierto. Al leer la figura, conviene confirmar primero la metric, el model size, la input resolution, el prompting y el test split de cada eje, y luego comparar cada rubro, en lugar de comparar solo el área que cubre el gráfico de radar.

\lecturefigure{slide-16-cogvlm-benchmarks.jpg}{Resumen de benchmarks de CogVLM al momento de la clase e información de código abierto}{Grabación oficial de Stanford Online 00:39:27}

\begin{warningbox}{Una benchmark table no equivale a production robustness}
VQA accuracy, caption score, grounding IoU y LLaVA-Bench miden problemas distintos; de ellos no se deducen directamente la GUI reliability, la cola larga del OCR, la tasa de alucinaciones, la adversarial robustness ni la tasa de éxito en tareas reales de usuarios. El número de descargas también indica solo adoption, no correctness.
\end{warningbox}

\subsection{CogAgent: por qué es necesaria una rama auxiliar de alta resolución}

Una screenshot de una página web contiene mucho texto pequeño, icons y relaciones de coordenadas. Si la resolución aumenta $s$ veces y el patch size no cambia, el número de tokens bidimensionales crece aproximadamente $s^2$ veces; el full attention cost puede acercarse incluso a $s^4$ veces. CogAgent usa una cross-attention high-resolution branch para que las características de alta resolución no tengan que expandirse todas al LLM hidden width.

\lecturefigure{slide-17-cogagent-architecture.jpg}{CogAgent: se agrega una high-resolution cross-attention branch para GUI/OCR}{Grabación oficial de Stanford Online 00:40:13}

Si los language tokens de baja resolución son $H_l\in\mathbb{R}^{n\times d}$ y las features de alta resolución son $H_h\in\mathbb{R}^{m\times d_h}$, la rama auxiliar puede escribirse como
\[
\widetilde H_l=H_l+\operatorname{CrossAttn}(H_lW_Q,H_hW_K,H_hW_V),
\]
cuyo objetivo es que las $m$ features de grano fino se consulten, en lugar de que todas se conviertan en LLM sequence tokens equivalentes.

\teachervoice{01:15:00--01:16:15, en la sesión de Q\&A el ponente distingue directamente CogAgent de CogVLM: el primero se obtiene del segundo mediante fine-tune/extensión, está orientado específicamente a la high-resolution web screen y usa una cross-attention más ligera para controlar el costo.}

\subsection{Web trace: el modelo produce una grounded action, no la frase «sabe usar el navegador»}

La tarea de ejemplo pide buscar el best paper de CVPR 2023. La slide muestra una trace de varios pasos: comprender la screenshot, ubicar el cuadro de búsqueda, escribir la query, leer los resultados y seguir haciendo clic. Una action suele requerir un type y coordinates/text, por ejemplo
\[
a_t=(\text{action\_type},x_t,y_t,\text{text}_t),
\qquad
p(a_t\mid I_t,g,h_{<t}).
\]

\lecturefigure{slide-18-cogagent-web-demo.jpg}{CogAgent web navigation: de la screenshot a la action sequence}{Grabación oficial de Stanford Online 00:41:19}

\begin{lstlisting}[caption={Action schema mínimo con capacidad de auditoría para un GUI agent}]
observation = capture_screenshot()
action = model.predict(goal, observation, history)
validate(action.type, action.coordinates, action.text)
execute_in_sandbox(action)
record(observation, action, result, failure_reason)
stop_only_when(goal_is_verified)
\end{lstlisting}

\begin{warningbox}{Una sola trace exitosa no estima la confiabilidad}
Un GUI agent real también requiere medir repeated trials, layout shift, small-text OCR, unexpected popups, permission boundary, irreversible actions y recovery. La demo de la clase demuestra que la interfaz es viable, no que la autonomous browsing ya sea robusta.
\end{warningbox}

\subsection{Vary: ampliar el vision vocabulary al servicio del OCR}

Vary combina distintas features visuales en un vision vocabulary más rico, con el objetivo de compensar las carencias de las CLIP features generales ante los detalles de document/OCR. Puede abstraerse como
\[
Z=\operatorname{Project}\!\left([f_{\text{general}}(I);f_{\text{ocr}}(I)]\right),
\]
donde los dos encoders tienen inductive biases distintos y, tras combinarse, se conectan al LLM.

\lecturefigure{slide-19-vary.jpg}{Vary: ensemble de vision features para ampliar el vocabulario visual}{Grabación oficial de Stanford Online 00:41:46}

\begin{knowledgebox}{El vision vocabulary no es un vocabulario discreto de palabras}
Aquí, vocabulary se refiere a los patrones de features visuales que pueden expresarse: layout, small text, objects, regions y estructura del documento. Agregar un encoder especializado puede mejorar el OCR, pero también aumenta los costos de entrenamiento, despliegue, feature alignment y mantenimiento.
\end{knowledgebox}

\subsection{GLM-4V: stride convolution simple y mixed training}

La estructura de GLM-4V que se mostró en clase vuelve a un LLaVA-style path más sencillo y solo sustituye el projector por una stride convolution para comprimir el feature map de alta resolución. Si el feature map de entrada es de $H\times W$ y el stride es $s$, el número de tokens baja aproximadamente de $HW$ a
\[
n_v\approx \frac{H}{s}\frac{W}{s}.
\]
Después, los datos de text/image se someten a mixed training para reducir el daño que el entrenamiento multimodal causa al language behavior.

\lecturefigure{slide-20-glm4v-architecture.jpg}{GLM-4V tal como se mostró en clase: stride convolution y text-image mixed training}{Grabación oficial de Stanford Online 00:42:34}

\teachervoice{00:41:46--00:43:00, el ponente señala deliberadamente que el modelo más reciente usa, de hecho, una architecture más simple. Este ejemplo respalda su data thesis: un bridge más complejo no produce automáticamente un sistema final más fuerte.}

\subsection{Tabla de resultados de GLM-4V: cada columna responde una pregunta distinta}

En la tabla aparecen a la vez MMBench, MME, MMMU, MathVista, MMVet, DocVQA, OCRBench, entre otros. Cubren capacidades distintas, como general perception, knowledge/reasoning, math/diagram y document OCR, por lo que no se puede calcular una «puntuación total» sin definición.

\lecturefigure{slide-21-glm4v-results.jpg}{Tabla de resultados de GLM-4V en varios benchmarks al momento de la clase}{Grabación oficial de Stanford Online 00:43:00}

\begin{knowledgebox}{Orden para leer una tabla de benchmarks multimodales}
Primero confirma el model/version y la resolución de entrada; después confirma zero-shot/few-shot, el prompt, el test split y la metric; luego compara por familia de tareas, en lugar de promediar directamente entre columnas. Si un modelo es fuerte en OCR pero débil en MMMU, la conclusión es que su distribución de capacidades es distinta, no simplemente «cuál es más inteligente».
\end{knowledgebox}

\begin{warningbox}{El GLM-4V de la clase no debe mezclarse con versiones posteriores}
Esta slide es anterior a los technical reports estables posteriores. Las notas solo registran la architecture y la table mostradas en ese momento, y no usan las capacidades de los GLM-4V/GLM-4.1V posteriores para demostrar retroactivamente las de la versión de 2024.
\end{warningbox}

\subsection{Resumen de la sección}

CogVLM usa experts para separar las rutas de parámetros, CogAgent usa high-resolution cross attention para controlar el costo de GUI/OCR, Vary amplía las features visuales y GLM-4V usa stride convolution para comprimir la entrada. Todos responden la misma pregunta: cómo llevar los detalles visuales al LLM y, al mismo tiempo, controlar el olvido, el número de tokens y el costo del sistema.
\section{Autoregressive image modeling: por qué una interfaz unificada no ganó en todas las tareas}

Si una image puede discretizarse en tokens, la idea más natural es seguir el camino de GPT: text-to-image, image-to-text e image completion son solo distintos ordenamientos de tokens y distintas mask. Esta línea tiene la elegancia de una interfaz unificada, pero la high-resolution sequence length, la tokenizer information loss y la sequential decoding la ponen frente a una fuerte competencia de diffusion en calidad de generación y en velocidad.

\subsection{CogView/DALL-E/Parti: primero convertir la imagen en una secuencia discreta}

Sea un image tokenizer que codifica la imagen $I$ en codes $z=(z_1,\ldots,z_m)$, y sean $x$ los tokens del texto. La factorization de text-to-image es
\[
p(z\mid x)=\prod_{i=1}^{m}p(z_i\mid x,z_{<i}).
\]
Para entrenar basta con concatenar los text tokens y los image tokens en una sola sequence, pero el número de image tokens puede llegar a miles, y la generación debe recorrer en serie todas las posiciones.

\lecturefigure{slide-22-autoregressive-text-to-image.jpg}{Autoregressive text-to-image: image tokenizer más GPT sequence}{Grabación oficial de Stanford Online 00:46:47}

\begin{importantbox}{El image tokenizer aporta a la vez compresión y pérdida}
El VQ tokenizer comprime los píxeles continuos en un codebook finito, lo que acorta mucho la secuencia; pero el texto fino, las texturas, la geometría y la precisión del color pueden perderse en la etapa de codificación. El Transformer posterior no puede recuperar información que el tokenizer nunca conservó.
\end{importantbox}

\teachervoice{00:45:00--00:49:32, el ponente coloca CogView y DALL-E/Parti en una misma recipe sencilla: primero discretizar y luego autoregress. La sencillez es una ventaja, y también el origen de los cuellos de botella en velocidad y en detalle.}

\subsection{Universal modeling: poder hacer todas las direcciones no equivale a ser el mejor en cada una}

Un mismo sequence model puede expresar varias tareas mediante ordenamientos y mask:
\[
\begin{aligned}
\text{text}\rightarrow\text{image}:&\ [x;z],\\
\text{image}\rightarrow\text{text}:&\ [z;x],\\
\text{masked image}:&\ [z_{\text{visible}};z_{\text{target}}].
\end{aligned}
\]
Esto sí unifica la API, pero en image understanding puede quedar por debajo de un VLM que conserva features continuas, y en image generation puede quedar por debajo de diffusion.

\lecturefigure{slide-23-universal-modeling.jpg}{Universal vision-language modeling: unificación por ordenamiento y compromiso de desempeño}{Grabación oficial de Stanford Online 00:49:32}

\begin{warningbox}{Un objective unificado y una solución óptima unificada son dos cosas distintas}
Que un modelo pueda expresar la likelihood de varias tareas no significa que, con parámetros, datos y compute limitados, alcance a la vez la Pareto frontier de cada uno de los specialized systems. Una representación compartida aporta transfer, pero también interference y bottleneck inadecuados.
\end{warningbox}

\teachervoice{00:49:32--00:52:11, el ponente hace una valoración muy mesurada de la línea de universal modeling en la que él mismo participó: logró la unificación, pero en image generation no supera a diffusion, y en image understanding tampoco supera a los VLM con características visuales continuas.}

\subsection{Resumen de la sección}

La mayor ventaja de autoregressive image modeling es la interfaz unificada y la madura language-model machinery; sus mayores problemas son la pérdida del tokenizer discreto, la generación en serie de secuencias largas y la linealización de las relaciones espaciales bidimensionales. No es que no pueda generar buenas imágenes, sino que la practical frontier se desplazó hacia diffusion, que se adapta mejor al cómputo paralelo de imágenes.
\section{Diffusion y Transformer: cambia la función objetivo, pero la experiencia de sistemas se sigue reutilizando}

Diffusion no predice la imagen siguiendo el orden de los tokens, sino que aprende un vector field de eliminación de ruido en varios noise levels. Así, cada paso procesa la latent completa, y la GPU utilization suele ser mejor que la de autoregressive decoding con lote de tamaño uno. Después, DiT volvió a cambiar el denoiser de U-Net a Transformer, lo que muestra que «el cambio de objective» no significa que Transformer quede fuera.

\subsection{DDPM: forward agrega ruido, reverse aprende a eliminarlo}

El autoregressive model de la sección anterior tiene que generar en serie a lo largo de la image-token sequence. En esta sección cambia el planteamiento del problema: ya no se pregunta «cuál es el siguiente token», sino «dado cierto noise level, cómo llevar la noisy latent completa hacia una dirección más limpia». Al leer las fórmulas hay que distinguir a la vez la forward corruption, definida por una persona, y el reverse denoising que aprende el modelo.

Se definen el noise schedule $\beta_t$, $\alpha_t=1-\beta_t$ y $\bar\alpha_t=\prod_{s=1}^{t}\alpha_s$. El forward process admite un muestreo en forma cerrada:
\[
q(x_t\mid x_0)=\mathcal{N}\!\left(\sqrt{\bar\alpha_t}x_0,(1-\bar\alpha_t)I\right),
\quad
x_t=\sqrt{\bar\alpha_t}x_0+\sqrt{1-\bar\alpha_t}\epsilon.
\]
El modelo suele predecir $\epsilon$:
\[
\mathcal{L}_{\text{simple}}=\mathbb{E}_{x_0,t,\epsilon}
\left[\|\epsilon-\epsilon_\theta(x_t,t,c)\|_2^2\right],
\]
donde $c$ es la text condition.

\lecturefigure{slide-24-diffusion-ddpm.jpg}{DDPM: noise schedule, denoising network y cálculo en paralelo de la imagen completa}{Video oficial de Stanford Online 00:52:11}

\begin{knowledgebox}{Por qué diffusion aprovecha más fácilmente toda la GPU}
Cada denoising step actualiza al mismo tiempo la latent completa, con matrices grandes y un alto grado de paralelismo. Cuando el lote es muy pequeño, un autoregressive image model produce un solo token por paso, y es difícil amortizar el kernel launch y el memory movement. Diffusion sigue necesitando muchos steps, pero el aprovechamiento del hardware en cada paso suele ser mejor.
\end{knowledgebox}

\teachervoice{01:12:30--01:15:00, en la sesión de Q\&A el ponente se niega a formular la conclusión como «imposible»: un autoregressive model lo bastante grande también puede generar imágenes de alta calidad. De lo que está más seguro es de la practical disadvantage que causan la sequential latency y las relaciones de larga distancia en dos dimensiones.}

\subsection{Relay Diffusion: al cambiar de resolución no basta con copiar el mismo ruido independiente}

Un mismo nominal noise level produce un grado de desenfoque visual distinto en cada resolution, porque la señal de las imágenes naturales tiene correlación espacial, mientras que el standard Gaussian noise suele ser independiente por pixel. Relay Diffusion usa correlated/block noise para ajustar la escala, de modo que las coarse-to-fine stages sean más coherentes en frecuencia y en SNR.

\lecturefigure{slide-25-relay-diffusion.jpg}{Relay Diffusion: desacoplar el noise schedule, la network y la resolution}{Video oficial de Stanford Online 00:54:49}

Si el block size es $k$, el noise de alta resolución puede escribirse como el sobremuestreo del noise de baja resolución más un residual:
\[
\epsilon_h=U_k(\epsilon_l)+\lambda\epsilon_{\text{res}},
\]
y con $\lambda$ se ajusta la variance de cada banda de frecuencia, para que la signal-to-noise ratio de la coarse stage y la de la fine stage sean más comparables. No se trata de un simple resize image, sino de redefinir el stochastic path entre resoluciones.

\begin{warningbox}{La «alineación de SNR» depende del schedule y del data spectrum}
La eficacia del block-correlated noise debe comprobarse con la resolution, el latent encoder, el noise schedule y los datos de entrenamiento concretos. Ofrece una explicación del mecanismo, pero no implica que toda coarse-to-fine diffusion pueda reutilizar directamente los mismos hiperparámetros.
\end{warningbox}

\teachervoice{00:52:11--00:55:16, el ponente enfatiza que, con el mismo noise independiente, la baja y la alta resolución no tienen la misma dificultad perceptual; lo esencial de Relay es desacoplar la resolution del schedule.}

\subsection{CogView3: el scale-up y la distillation forman parte del sistema}

Relay Diffusion resuelve el schedule entre resoluciones. Para convertir este mecanismo en un text-to-image system utilizable, también hay que ampliar al mismo tiempo el modelo y los datos, y comprimir el costoso teacher sampling en menos steps. CogView3 es precisamente esa extensión sistemática, y en la slide informa la velocidad de generación de imágenes de $1024\!\times\!1024$ del distilled model. Toda cifra de velocidad debe ir ligada al hardware, la implementación, el sampling step, el lote y la configuración de calidad.

\lecturefigure{slide-26-cogview3.jpg}{CogView3: ampliar Relay Diffusion y acelerarlo mediante distillation}{Video oficial de Stanford Online 00:55:16}

\begin{importantbox}{El objetivo de la distillation no es solo reducir parámetros}
La diffusion distillation suele comprimir una teacher trajectory de muchos pasos en menos sampling steps. Lo que realmente se optimiza es el tradeoff entre end-to-end latency y quality; aun sin cambiar la cantidad de parámetros, puede acelerarse mucho al reducir las function evaluations.
\end{importantbox}

\subsection{DiT: sustituir el denoiser por un Transformer}

DiT trata los latent patches como tokens e incorpora la timestep/class condition mediante adaptive LayerNorm. Si la normalization ordinaria es $\operatorname{LN}(h)$, adaLN puede escribirse como
\[
\operatorname{adaLN}(h;c)=\gamma(c)\odot\operatorname{LN}(h)+\beta(c),
\]
donde $c$ se obtiene del timestep y del embedding de la condición. Así, cada capa puede ajustar la feature scale y el shift según el noise level.

\lecturefigure{slide-27-diffusion-transformer.jpg}{DiT: Transformer blocks y adaptive normalization en lugar de U-Net}{Video oficial de Stanford Online 00:57:02}

\begin{knowledgebox}{Diferencia entre AdaLN y cross attention}
Con cross attention, los image tokens consultan un conjunto de condition tokens; adaLN, en cambio, comprime la condition en scale/shift/gate para cada capa. La primera conserva la alineación a nivel de token; la segunda se parece más a una modulación global. Un sistema real puede usar ambas a la vez.
\end{knowledgebox}

\teachervoice{00:55:16--00:57:43, el ponente observa que la entrada del timestep es solo una condition muy pequeña y, sin embargo, puede pasar por un MLP grande para generar los parámetros de modulación de cada capa; lo considera un diseño de ingeniería que todavía admite simplificación.}

\subsection{Stable Diffusion 3 y MM-DiT: reaparecen los text/image experts}

MM-DiT mantiene parameter paths separados para los text tokens y los image tokens, y luego los hace interactuar mediante attention. Puede abstraerse como
\[
Q_t=H_tW_Q^{(t)},\quad Q_i=H_iW_Q^{(i)},
\]
mientras que la joint attention intercambia información sobre $K,V$ concatenados. La idea que comparte con CogVLM son las modality-specific transformations; la diferencia es que la tarea es diffusion denoising y no autoregressive response.

\lecturefigure{slide-28-stable-diffusion3-mmdit.jpg}{Stable Diffusion 3: los text/image expert paths de MM-DiT}{Video oficial de Stanford Online 00:57:43}

\begin{warningbox}{Nombres de módulo parecidos no implican el mismo problema de entrenamiento}
Los vision experts de CogVLM sirven al image-conditioned language modeling; los modality paths de MM-DiT sirven al rectified flow/diffusion denoising. Separar los parámetros es un patrón de diseño compartido, pero la loss, la sequence, la condition y la evaluation son completamente distintas.
\end{warningbox}

\subsection{Resumen de la sección}

Diffusion se impuso en la practical image generation, sobre todo gracias al paralelismo sobre la imagen completa, a la latent continua y a una vía de modelado espacial más adecuada. Relay se ocupa del noise entre resoluciones, la distillation reduce los sampling steps, y DiT/MM-DiT devuelven el Transformer al denoiser. La función objetivo y la vía de sistemas determinan juntas la eficiencia final.
\section{Video generation y agenda de investigación: formular las predicciones como preguntas verificables}

El video acumula todos los cuellos de botella anteriores: más token/latent, temporal consistency más difícil, high-resolution más costosa, y data coverage y caption quality peores. El ponente no presentó conclusiones internas sobre Sora; a partir de los fenómenos públicos, desglosó un conjunto de hipótesis de ingeniería de 2024 y, con base en ellas, propuso líneas de investigación.

\subsection{Sora-like recipe: 3D latents, scale, context y data}

La clase toma CogVideo como referencia temprana de text-to-video autoregressive y luego divide la Sora-like improvement en deflickering, image quality, resolution, long-context infrastructure y recaption/coverage. Si el latent del video tiene tamaño $T\times H\times W$, el número naive de tokens es
\[
n_v=T\frac{H}{p_h}\frac{W}{p_w},
\]
y el costo cuadrático de la full attention se dispara rápidamente; por ello, la 3D latent compression, la factorized attention, el parallelism y el data pipeline son condiciones necesarias.

\lecturefigure{slide-29-video-generation.jpg}{Video generation: de CogVideo a la Sora-like scaling recipe}{Grabación oficial de Stanford Online 00:58:29}

\begin{knowledgebox}{Qué resuelve cada uno de los cinco factores de ingeniería}
\small\sloppy
\begin{tabularx}{\textwidth}{L{0.19\textwidth} L{0.32\textwidth} X}
\toprule
Factor & Problema principal & Capacidad que no garantiza por sí solo \\
\midrule
3D latent encoder & Comprimir el espacio-tiempo y reducir el flicker & Física real y causalidad de largo plazo \\
Model/data scale & Mejorar la calidad visual y la cobertura & Ausencia de sesgo, seguridad, controlabilidad \\
High resolution & Conservar texto y detalles & Movimiento y geometría correctos \\
Long-context infra & Soportar más frames/tokens & Aprovechar efectivamente todo el context \\
Recaption y coverage & Mejorar la text-video supervision & Aprender leyes físicas a partir de raw video \\
\bottomrule
\end{tabularx}
\end{knowledgebox}

\teachervoice{00:57:54--01:04:00, el ponente califica estos factores como un resumen del Sora-like progress, no como una controlled ablation. Por eso, las notas los presentan como hypotheses verificables y no como una receta interna conocida.}

\subsection{Predicciones de 2024: reconocimiento, comprensión de video y embodied AI}

La sección anterior desglosó cómo los modelos y los datos amplían la video generation; aquí se pasa a un research forecast más incierto. La slide de predicciones sostiene que el reconocimiento visual de alto nivel se abaratará, que la video understanding será un nuevo tema candente y que la embodied AI producirá demos llamativas, pero difícilmente entrará en la vida cotidiana a corto plazo. La forma correcta de registrarlo es conservar la fecha (mayo de 2024), los supuestos de recursos y el tono de «posiblemente», en lugar de seleccionar a posteriori solo las frases que acertaron.

\lecturefigure{slide-30-research-predictions.jpg}{Predicciones de la clase en 2024: reconocimiento visual, video understanding y embodied AI}{Grabación oficial de Stanford Online 01:04:00}

\begin{warningbox}{Una prediction no es un benchmark result}
«Uno o dos años», «80\% solved» y «no podrá influir pronto en la vida real» no son criterios medibles directamente. Para evaluar una predicción hay que redefinir la task distribution, el cost, la accuracy, el long-tail, el deployment environment y el momento de la medición.
\end{warningbox}

\teachervoice{01:00:00--01:04:00, el ponente expresa a la vez optimismo y límites: el reconocimiento visual se abaratará con rapidez, pero los video models todavía tienen problemas de hallucination/counting; la embodied AI tendrá demos sorprendentes, pero podría ser costosa y difícil de desplegar en la vida diaria.}

\subsection{Recomendaciones de investigación: elegir el problema según recursos y objetivos}

La última página no es una «clasificación de temas», sino una forma de relacionar el career objective, el compute access y el research risk. Los video datasets/benchmarks son adecuados para equipos con menos recursos pero capaces de organizar la evaluation; se considera que speech tiene alta demanda pero poca inversión; architecture/optimizer puede tener gran impacto, pero exige capacidades de hardware y de sistemas; y compute-to-data es la dirección central una vez que el web corpus se va saturando.

\lecturefigure{slide-31-research-advice.jpg}{Elegir problemas de investigación según objetivos y recursos: video, speech, systems, architecture y data}{Grabación oficial de Stanford Online 01:07:34}

\begin{knowledgebox}{Reescribir las recomendaciones como matriz de decisión}
\small\sloppy
\begin{tabularx}{\textwidth}{L{0.17\textwidth} L{0.20\textwidth} L{0.20\textwidth} X}
\toprule
Dirección & Activo principal & Riesgo principal & Entregable verificable \\
\midrule
Datos/evaluación de video & collection, annotation, evaluation & Fugas, shortcut, cobertura insuficiente & dataset card, held-out test, error taxonomy \\
Speech/audio & Escenarios de usuario y datos multilingües & Privacidad, ruido, latencia en tiempo real & Métricas de ASR, TTS y de diálogo de extremo a extremo \\
Colaboración con sistemas & Topología de GPU y conocimiento de kernels & Reproducibilidad y dependencia del hardware & Throughput, latencia, utilización y costo \\
Arquitectura/optimizador & controlled runs a gran escala & Cómputo costoso, conclusiones inestables & matched-compute scaling/ablation \\
Compute-to-data & verifier, execution, search & feedback bias, synthetic collapse & provenance, verification rate, novelty \\
\bottomrule
\end{tabularx}
\end{knowledgebox}

\begin{lstlisting}[caption={Verificación previa a la publicación de un proyecto de investigación multimodal}]
freeze_task_definition_and_data_split()
record_model_version_resolution_and_prompt()
report_compute_latency_memory_and_cost()
separate_demo_success_from_repeated_trial_rate()
audit_shortcuts_leakage_and_long_tail_failures()
publish_provenance_metrics_and_negative_results()
\end{lstlisting}

\teachervoice{01:05:00--01:08:10, el ponente afirma que la speech AI está subestimada y recomienda a quienes trabajan en algoritmos colaborar cuanto antes con systems researchers, porque los mejores algoritmos deben adaptarse al hardware contemporáneo.}

\subsection{Resumen de la sección}

La video generation no consiste en ejecutar un image model durante más frames, sino en un problema de sistemas que integran la compresión espaciotemporal, el entrenamiento en paralelo, el long context, la cobertura de datos y la evaluación. Las recomendaciones de investigación también deben entenderse en función de los recursos y de los entregables verificables, en lugar de tomar las predicciones de la clase como una hoja de ruta segura.
\section{Ampliación de la sesión de Q\&A: limitaciones que las slides de la clase no mencionan}

La sesión de Q\&A no añadió ningún visual state independiente, pero completó varias acotaciones que determinan si las conclusiones técnicas son confiables: el prefill cost del long context, la relación de suma no nula entre data y architecture, la incertidumbre sobre la ventaja de diffusion, el origen de la alta resolución de CogAgent y si un modelo de video puede aprender física a partir de la raw signal.

\subsection{Long context, data y architecture: tres equilibrios}

Primero, un contexto largo puede reducir la necesidad de un retrieval pipeline externo, pero la prefill latency y la interferencia de tokens irrelevantes siguen presentes. Segundo, los data pueden expresar muchos task-specific biases, pero una general architecture improvement todavía puede cambiar la compute frontier de todas las tareas. Tercero, colocar todos los documents o reasoning branches en el context exige que el modelo aprenda a aprovecharlos; no basta con ampliar la ventana.

\begin{importantbox}{Conclusiones de ingeniería de la sesión de Q\&A}
\begin{enumerate}
  \item Para documentos largos con respuestas cortas, puede ser razonable aceptar un prefill más lento; para un agent en tiempo real, no necesariamente.
  \item Ante un comportamiento específico, conviene intentar primero con data/evaluation; ante un cuello de botella que se repite entre tareas, considerar después la architecture.
  \item Si el reasoning procedure incluye rutas fallidas, colocarlas explícitamente en el context puede aportar más información que conservar solo el beam score, aunque esto todavía requiere experimentos.
\end{enumerate}
\end{importantbox}

\teachervoice{01:18:10--01:19:59, el ponente se inclina por colocar las alternative paths junto con la failure information en el context, para que el modelo aprenda a compararlas, en lugar de usar solo una hard-coded beam probability; es una observación empírica, no un universal theorem.}

\subsection{La video understanding no adquiere automáticamente el mundo físico}

Los VLM/video models actuales dependen sobre todo de text-image o text-video pairs. Si la training signal proviene principalmente de annotation humana, el modelo aprende los fenómenos que cubren las descripciones humanas, y no necesariamente recupera a partir del raw motion la conservation, el contact, la object permanence o las causal interventions.

\begin{warningbox}{«Haber visto muchos videos» no equivale a «haber aprendido física»}
Para aprender física a partir de video sin anotar pueden hacer falta temporal prediction, masked video modeling, world-model objectives, interaction data o una simulation verificable. Ampliar solo los annotated pairs todavía puede llevar a aprender shortcuts del lenguaje y correlaciones superficiales.
\end{warningbox}

\teachervoice{01:16:15--01:18:10, el ponente afirma con claridad que, si la fuente de datos no contiene physical rules, el preentrenamiento actual basado en pairs no las aprenderá de forma automática; considera que la nueva video self-supervision es un problema de investigación necesario.}

\subsection{Resumen de la sección}

La sesión de Q\&A convierte los headlines de la exposición principal en juicios de ingeniería con condiciones: el context tiene latencia, los data no eliminan la architecture, diffusion no ha superado matemáticamente a autoregression, una GUI demo no representa un agent confiable y la video quantity tampoco equivale a physical understanding.
\section{Síntesis y ampliación}

Lo que más vale conservar de esta clase no es la lista de modelos, sino un marco de análisis de tres niveles. El primer nivel es \textbf{representation/objective}: qué información conserva cada uno de language token, continuous vision feature, discrete image code y noisy latent. El segundo nivel es \textbf{architecture/systems}: cómo projection, experts, cross attention, ZeRO, parallelism y long context llevan la representación a un hardware en el que se puede entrenar. El tercer nivel es \textbf{data/evaluation}: cómo instruction, preference, caption, GUI trace, video coverage y benchmark definen el comportamiento que el modelo termina aprendiendo.

\subsection{Unificación de los tres ejes}

\begin{importantbox}{Panorama unificado final del curso}
\begin{enumerate}
  \item \textbf{El LLM aporta la base secuencial de uso general: }scaling, alignment y systems hacen del modelo grande una interfaz reutilizable.
  \item \textbf{La multimodality expone el cuello de botella de la interfaz: }los detalles visuales, la alta resolución, el espacio bidimensional y la estructura temporal del video obligan al modelo a cambiar el bridge, el objective y la compression.
  \item \textbf{Los data determinan los límites del comportamiento: }la architecture puede ampliar el rango que el modelo es capaz de ajustar, pero benchmark, annotation, recaption, verification y failure coverage determinan lo que el sistema aprende en realidad.
\end{enumerate}
\end{importantbox}

\begin{warningbox}{Los tres errores de generalización más frecuentes}
No conviertas «la misma loss de preentrenamiento» en un teorema de «la misma capacidad»; no conviertas una sola GUI trace en agent reliability; no conviertas la practical advantage de diffusion en una prueba de imposibilidad de la autoregression. Los tres errores provienen de borrar las diferencias entre los niveles de evidencia presentados en clase.
\end{warningbox}

\subsection{Checklist para reproducir experimentos}

Para reproducir cualquiera de los modelos de esta clase, se debe publicar al mismo tiempo, como mínimo, architecture version, data mixture/provenance, resolution/tokenización, training compute, inference steps, prompt/evaluation protocol y repeated-trial failures. Publicar solo el headline score impide distinguir entre data, algorithm y architecture.

\begin{lstlisting}[caption={Lista unificada de reproducción para LLM/VLM/diffusion}]
pin(code_commit, model_checkpoint, tokenizer_or_vae)
document(data_sources, filters, synthetic_fraction, licenses)
report(train_flops, gpu_hours, parallelism, peak_memory)
report(input_resolution, context_length, sampling_steps)
freeze(prompts, decoding, benchmark_version, test_split)
publish(ablations, confidence_intervals, failures, costs)
\end{lstlisting}

\subsection{Preguntas de autoevaluación}

\begin{multicols}{2}
\footnotesize
\begin{enumerate}
  \setlength{\itemsep}{0pt}
  \setlength{\parskip}{0pt}
  \item ¿En qué difiere la información de condicionamiento de GLM blank infilling respecto de un MLM puro y de una autoregression pura?
  \item ¿Por qué la scaling law es un asignador de presupuesto y no una garantía de rendimiento?
  \item ¿Por qué la perplexity no puede compararse directamente entre tokenizers distintos?
  \item ¿Qué tipo de inference memory reduce principalmente GQA?
  \item ¿A qué aplica sharding cada uno de ZeRO-1/2/3?
  \item ¿En qué se diferencia activation checkpointing de model checkpointing?
  \item ¿Qué dimensión divide cada uno de TP, PP y context parallelism?
  \item ¿Por qué el principal costo de interacción de long context puede estar en el prefill?
  \item DPO no tiene un reward model independiente; ¿por qué sigue dependiendo de reward assumptions?
  \item ¿Cómo muestra el ejemplo de CogQA que data/algorithm/architecture son intercambiables?
  \item ¿Qué dos tipos de trabajo realiza a la vez el Q-Former de BLIP-2?
  \item La projection de LLaVA es sencilla, pero ¿por qué el OCR todavía puede ser deficiente?
  \item ¿Qué capacidad buscan proteger los vision experts de CogVLM?
  \item ¿Por qué CogAgent necesita una ruta lateral de alta resolución?
  \item ¿Qué capacidades no puede demostrar una sola web trace?
  \item ¿En qué difieren Vary y GLM-4V en el tratamiento del OCR bottleneck?
  \item ¿Qué ventajas y desventajas aporta el image tokenizer a la autoregressive generation?
  \item ¿Por qué diffusion suele aprovechar mejor la GPU que la AR decoding con lote de tamaño uno?
  \item ¿Por qué Relay Diffusion necesita cambiar la noise correlation entre resolutions?
  \item ¿En qué difieren las rutas de condition de adaLN y de cross attention?
  \item En una Sora-like recipe, ¿qué es un problema de representación y qué es un problema de systems/data?
  \item ¿Por qué los annotated video pairs no garantizan que se aprendan las physical rules?
\end{enumerate}
\end{multicols}

\subsection{Lecturas adicionales}

\begin{multicols}{2}
\footnotesize
\begin{itemize}
  \setlength{\itemsep}{0.15em}
  \setlength{\parskip}{0pt}
  \item Du et al., \href{https://arxiv.org/abs/2103.10360}{GLM: General Language Model Pretraining with Autoregressive Blank Infilling}: para entender la unificación de las funciones objetivo.
  \item Rajbhandari et al., \href{https://arxiv.org/abs/1910.02054}{ZeRO}; Shoeybi et al., \href{https://arxiv.org/abs/1909.08053}{Megatron-LM}: para entender la partición del estado y el paralelismo de modelo.
  \item Rafailov et al., \href{https://arxiv.org/abs/2305.18290}{Direct Preference Optimization}: para entender el pairwise preference objective.
  \item Li et al., \href{https://arxiv.org/abs/2301.12597}{BLIP-2}; Liu et al., \href{https://arxiv.org/abs/2304.08485}{LLaVA}: para comparar el Q-Former con el projection bridge.
  \item Wang et al., \href{https://arxiv.org/abs/2311.03079}{CogVLM}; Hong et al., \href{https://arxiv.org/abs/2312.08914}{CogAgent}: para entender los vision experts y el GUI high-resolution path.
  \item Ho et al., \href{https://arxiv.org/abs/2006.11239}{DDPM}; Peebles and Xie, \href{https://arxiv.org/abs/2212.09748}{DiT}; Esser et al., \href{https://arxiv.org/abs/2403.03206}{Stable Diffusion 3}: del diffusion objective al Transformer denoiser.
  \item Snell et al., \href{https://arxiv.org/abs/2409.01236}{Same Pre-training Loss, Better Downstream}: lectura post-lecture que sirve de contraejemplo a la afirmación categórica de la clase de que «la loss determina la capacidad».
\end{itemize}
\end{multicols}

\end{document}
